% ---- score2 standard preamble (from: python3 score2_tools.py preamble) ----
\documentclass[10pt]{extarticle}
\usepackage[a4paper,margin=1.5cm]{geometry}
\usepackage[T1]{fontenc}
\usepackage{amsmath,amssymb}
\usepackage{graphicx}
\usepackage[export]{adjustbox}
\usepackage{multicol}
\usepackage{enumitem}
\usepackage{booktabs,array,tabularx}
\setlength{\columnsep}{0.7cm}
\setlength{\parindent}{0pt}
\setlength{\parskip}{3pt}
\pagestyle{empty}
\makeatletter
% figure registry (figdefs.tex is generated by the extract step)
\newcommand{\qfigdef}[4]{\@namedef{qf@file@#1}{#2}\@namedef{qf@w@#1}{#3}\@namedef{qf@wide@#1}{#4}}
\newcommand{\qf@check}[1]{\@ifundefined{qf@file@#1}{\PackageError{score2}{Unknown figure id #1}{}}{}}
\newcommand{\qf@img}[4][]{\includegraphics[width=#2\dimexpr\@nameuse{qf@w@#4}pt\relax,max width=#3,max height=0.45\textheight,#1]{\@nameuse{qf@file@#4}}}
% Full-width questions: finalize wraps every question that contains a wide figure in \qfullbegin...\qfullend,
% so the WHOLE question (stem, figures, options) is set across the page in reading order, then the two columns
% resume. The question's head - its number, stem and first figure (\qfullsplit marks the end of the head) -
% is kept on one page, so a number is never left at a page foot with its figure overleaf; the rest may break
% across pages normally, so no big gaps are left. Without a \qfullsplit, the whole question is kept together
% if it is under 45% of a page. Never write these yourself.
\newif\ifqf@full
\newif\ifqf@boxed
\newsavebox{\qf@box}
\newcommand{\qf@open}{\global\qf@fulltrue\par\global\qf@boxedtrue
  \setbox\qf@box\vbox\bgroup\hsize\textwidth\linewidth\textwidth}
\newcommand{\qf@emit}[1]{\par\egroup\global\qf@boxedfalse
  \ifdim\dimexpr\ht\qf@box+\dp\qf@box\relax<#1\textheight\noindent\usebox{\qf@box}\par\else\unvbox\qf@box\fi}
\newcommand{\qfullsplit}{\ifqf@boxed\qf@emit{0.85}\fi}
\newcommand{\qf@close}{\ifqf@boxed\qf@emit{0.45}\fi\global\qf@fullfalse}
\newcommand{\qfullbegin}{\end{multicols}\qf@open}
\newcommand{\qfullend}{\qf@close\begin{multicols}{2}}
\newcommand{\qfullstart}{\qf@open}
\newcommand{\qfulldone}{\qf@close}
% \qfig[scale]{id}: block figure, centred, directly where it is written (size set by the tool).
% Fallback if a wide figure is not inside a full-width question: break out of the columns in place - no floats.
\newcommand{\qfig}[2][1]{\qf@check{#2}%
  \ifqf@full\qf@blk{#1}{#2}\else
  \ifnum\@nameuse{qf@wide@#2}=1 \expandafter\@firstoftwo\else\expandafter\@secondoftwo\fi
  {\end{multicols}\begin{center}\qf@img{#1}{\textwidth}{#2}\end{center}\begin{multicols}{2}}%
  {\qf@blk{#1}{#2}}\fi}
\newcommand{\qf@blk}[2]{\par\smallskip{\centering\qf@img{#1}{\linewidth}{#2}\par}\smallskip}
% \qfigi[scale]{id}: inline figure (inside option lists / tables / minipages); never breaks columns.
\newcommand{\qfigi}[2][1]{\qf@check{#2}\qf@img[valign=c]{#1}{\linewidth}{#2}}
\makeatother
% structure macros - lint reads these to check completeness and figure placement
\newcommand{\papersec}[1]{\par\bigskip{\centering\Large\bfseries PAPER-#1\par}\medskip}
\newcommand{\partsec}[2]{\par\medskip\noindent{\large\bfseries PART-#1: #2}\par\smallskip}
\newcommand{\sect}[1]{\par\smallskip\noindent\textbf{#1}\par}
\newcommand{\qn}[1]{\par\smallskip\noindent\textbf{#1.}~}
\newcommand{\paraq}[2]{\par\smallskip\noindent\textbf{Paragraph for Questions #1 and #2}\par}
\newenvironment{qoptions}{\begin{itemize}[leftmargin=1.8em,itemsep=1pt,topsep=2pt,label={}]}{\end{itemize}}
\newcommand{\opt}[2]{\item[(#1)] #2}
% figure registry (generated by score2_tools.py finalize)
\qfigdef{p1_phy_s1_q04_1}{Score-II - Test 19 (20 Feb 2024)-img/p1_phy_s1_q04_1.png}{70.6}{0}
\qfigdef{p1_phy_s1_q05_1}{Score-II - Test 19 (20 Feb 2024)-img/p1_phy_s1_q05_1.png}{84.2}{0}
\qfigdef{p1_phy_s1_q06_1}{Score-II - Test 19 (20 Feb 2024)-img/p1_phy_s1_q06_1.png}{110.9}{0}
\qfigdef{p1_phy_s1_q11_1}{Score-II - Test 19 (20 Feb 2024)-img/p1_phy_s1_q11_1.png}{87.0}{0}
\qfigdef{p1_phy_s1_q13_1}{Score-II - Test 19 (20 Feb 2024)-img/p1_phy_s1_q13_1.png}{83.3}{0}
\qfigdef{p1_phy_s1_q16_1}{Score-II - Test 19 (20 Feb 2024)-img/p1_phy_s1_q16_1.png}{76.1}{0}
\qfigdef{p1_phy_s2_q01_1}{Score-II - Test 19 (20 Feb 2024)-img/p1_phy_s2_q01_1.png}{67.7}{0}
\qfigdef{p1_phy_s2_q04_1}{Score-II - Test 19 (20 Feb 2024)-img/p1_phy_s2_q04_1.png}{87.9}{0}
\qfigdef{p1_phy_s2_q08_1}{Score-II - Test 19 (20 Feb 2024)-img/p1_phy_s2_q08_1.png}{100.0}{0}
\qfigdef{p1_chem_s1_q15_1}{Score-II - Test 19 (20 Feb 2024)-img/p1_chem_s1_q15_1.png}{45.9}{0}
\qfigdef{p1_chem_s1_q15_2}{Score-II - Test 19 (20 Feb 2024)-img/p1_chem_s1_q15_2.png}{40.9}{0}
\qfigdef{p1_chem_s1_q15_5}{Score-II - Test 19 (20 Feb 2024)-img/p1_chem_s1_q15_5.png}{39.2}{0}
\qfigdef{p1_chem_s1_q15_3}{Score-II - Test 19 (20 Feb 2024)-img/p1_chem_s1_q15_3.png}{40.9}{0}
\qfigdef{p1_chem_s1_q15_4}{Score-II - Test 19 (20 Feb 2024)-img/p1_chem_s1_q15_4.png}{154.2}{0}
\qfigdef{p1_chem_s1_q16_1}{Score-II - Test 19 (20 Feb 2024)-img/p1_chem_s1_q16_1.png}{40.9}{0}
\qfigdef{p1_chem_s1_q16_2}{Score-II - Test 19 (20 Feb 2024)-img/p1_chem_s1_q16_2.png}{72.4}{0}
\qfigdef{p1_chem_s1_q16_3}{Score-II - Test 19 (20 Feb 2024)-img/p1_chem_s1_q16_3.png}{111.0}{0}
\qfigdef{p1_chem_s1_q16_4}{Score-II - Test 19 (20 Feb 2024)-img/p1_chem_s1_q16_4.png}{111.0}{0}
\qfigdef{p1_chem_s1_q18_1}{Score-II - Test 19 (20 Feb 2024)-img/p1_chem_s1_q18_1.png}{108.4}{0}
\qfigdef{p1_chem_s1_q18_2}{Score-II - Test 19 (20 Feb 2024)-img/p1_chem_s1_q18_2.png}{46.3}{0}
\qfigdef{p1_chem_s1_q18_3}{Score-II - Test 19 (20 Feb 2024)-img/p1_chem_s1_q18_3.png}{38.1}{0}
\qfigdef{p1_chem_s1_q18_4}{Score-II - Test 19 (20 Feb 2024)-img/p1_chem_s1_q18_4.png}{46.3}{0}
\qfigdef{p1_chem_s1_q18_5}{Score-II - Test 19 (20 Feb 2024)-img/p1_chem_s1_q18_5.png}{46.3}{0}
\qfigdef{p1_chem_s1_q19_1}{Score-II - Test 19 (20 Feb 2024)-img/p1_chem_s1_q19_1.png}{97.7}{0}
\qfigdef{p1_chem_s1_q19_2}{Score-II - Test 19 (20 Feb 2024)-img/p1_chem_s1_q19_2.png}{27.3}{0}
\qfigdef{p1_chem_s1_q19_3}{Score-II - Test 19 (20 Feb 2024)-img/p1_chem_s1_q19_3.png}{27.5}{0}
\qfigdef{p1_chem_s1_q19_4}{Score-II - Test 19 (20 Feb 2024)-img/p1_chem_s1_q19_4.png}{25.8}{0}
\qfigdef{p1_chem_s1_q19_5}{Score-II - Test 19 (20 Feb 2024)-img/p1_chem_s1_q19_5.png}{24.3}{0}
\qfigdef{p1_chem_s1_q20_2}{Score-II - Test 19 (20 Feb 2024)-img/p1_chem_s1_q20_2.png}{57.0}{0}
\qfigdef{p1_chem_s2_q08_1}{Score-II - Test 19 (20 Feb 2024)-img/p1_chem_s2_q08_1.png}{83.9}{0}
\qfigdef{p1_chem_s2_q09_1}{Score-II - Test 19 (20 Feb 2024)-img/p1_chem_s2_q09_1.png}{40.7}{0}
\qfigdef{p1_chem_s2_q09_2}{Score-II - Test 19 (20 Feb 2024)-img/p1_chem_s2_q09_2.png}{46.3}{0}
\qfigdef{p1_chem_s2_q09_3}{Score-II - Test 19 (20 Feb 2024)-img/p1_chem_s2_q09_3.png}{38.7}{0}
\qfigdef{p1_chem_s2_q09_4}{Score-II - Test 19 (20 Feb 2024)-img/p1_chem_s2_q09_4.png}{35.5}{0}
\qfigdef{p1_chem_s2_q09_7}{Score-II - Test 19 (20 Feb 2024)-img/p1_chem_s2_q09_7.png}{46.3}{0}
\qfigdef{p1_chem_s2_q09_6}{Score-II - Test 19 (20 Feb 2024)-img/p1_chem_s2_q09_6.png}{46.3}{0}
\qfigdef{p1_chem_s2_q09_5}{Score-II - Test 19 (20 Feb 2024)-img/p1_chem_s2_q09_5.png}{57.3}{0}
\qfigdef{p1_chem_s2_q10_1}{Score-II - Test 19 (20 Feb 2024)-img/p1_chem_s2_q10_1.png}{30.9}{0}
% ---- body skeleton ----
% \begin{document}
% \begin{center}{\Large\bfseries JEE (Main+Advanced) Score-II --- Test NN (DD Mon YYYY)}\end{center}
% \begin{multicols}{2}
%   \papersec{1}   (only for two-paper tests)
%   \partsec{1}{PHYSICS} \sect{Section-I (i): Single Correct} \qn{1} ... \qfig{p1_phy_q01_1} ...
% \end{multicols}
% \end{document}

\begin{document}
\begin{center}{\Large\bfseries JEE (Main+Advanced) Score-II --- Test 19 (20 Feb 2024)}\end{center}
\begin{multicols}{2}
\raggedcolumns
\partsec{1}{PHYSICS}
\sect{Section-I}

\qn{1} The resistance of a moving coil galvanometer is 20$\Omega$. It requires 0.01A current for full scale deflection. Calculate the value of resistance required to convert it into a voltmeter of range 20 volt :-
\begin{qoptions}
\opt{A}{1980 $\Omega$}
\opt{B}{2000 $\Omega$}
\opt{C}{1880 $\Omega$}
\opt{D}{1500 $\Omega$}
\end{qoptions}

\qn{2} Two long parallel wires carrying equal currents in opposite directions are placed at $x=\pm a$ parallel to y-axis with $z=0$. Magnetic field at origin O is $B_1$ and at P $(2a,0,0)$ is $B_2$. Then $B_1/B_2$ is :-
\begin{qoptions}
\opt{A}{3}
\opt{B}{$\dfrac12$}
\opt{C}{$\dfrac13$}
\opt{D}{2}
\end{qoptions}

\qn{3} A series R--L--C circuit consists of an 8.0 $\Omega$ resistor, a 5.0 $\mu$F capacitor and a 50.0 mH inductor. A variable frequency source applies an emf of 400 V (rms) across the combination. The power delivered to the circuit when the frequency is equal to one-half the resonance frequency is :-
\begin{qoptions}
\opt{A}{52 W}
\opt{B}{57 W}
\opt{C}{63 W}
\opt{D}{69 W}
\end{qoptions}

\qn{4} A ring of radius $R$ with a uniformly distributed charge $q$ as shown in figure. A change $q_0$ is now placed at the centre of the ring. Find the increment in the tension in ring :-
\par\nobreak
\qfig{p1_phy_s1_q04_1}
\begin{qoptions}
\opt{A}{$\dfrac{kqq_0}{2\pi R^2}$}
\opt{B}{$\dfrac{kaq_0}{R^2}$}
\opt{C}{$\dfrac{kqq_0}{\pi R^2}$}
\opt{D}{None of these}
\end{qoptions}

\qn{5} In the adjoining circuit, initially the switch S is open. The switch `S' is closed at $t=0$. The difference between the maximum and minimum current that can flow in the circuit through battery is : --
\qfig{p1_phy_s1_q05_1}
\begin{qoptions}
\opt{A}{2 Amp}
\opt{B}{3 Amp}
\opt{C}{1 Amp}
\opt{D}{Nothing can be concluded}
\end{qoptions}

\qn{6} In the given circuit if point C is connected to the earth and a potential of $+2000$ V is given to point A, then potential at B is :
\qfig{p1_phy_s1_q06_1}
\begin{qoptions}
\opt{A}{1500 V}
\opt{B}{1000 V}
\opt{C}{500 V}
\opt{D}{400 V}
\end{qoptions}

\qn{7} 5.6 liter of helium gas at STP is adiabatically compressed to 0.7 liter. Taking the initial temperature to be $T_1$, the work done in the process is :
\begin{qoptions}
\opt{A}{$\dfrac98RT_1$}
\opt{B}{$\dfrac32RT_1$}
\opt{C}{$\dfrac{15}{8}RT_1$}
\opt{D}{$\dfrac92RT_1$}
\end{qoptions}

\qn{8} A vessel contains a mixture of 1 mole of oxygen and two moles of nitrogen at 300 K. The ratio of the rotational kinetic energy per O$_2$ molecule to that per N$_2$ molecule is :
\begin{qoptions}
\opt{A}{1 : 1}
\opt{B}{1 : 2}
\opt{C}{2 : 1}
\opt{D}{depends on the moment of inertia of the two molecules}
\end{qoptions}

\qn{9} In a Young's double slit experiment the intensity at a point where the path different is $\dfrac{\lambda}{6}$ ($\lambda$ being the wavelength of the light used) is $I$. If $I_0$ denotes the maximum intensity, $I/I_0$ is equal to :-
\begin{qoptions}
\opt{A}{$\dfrac{1}{\sqrt2}$}
\opt{B}{$\dfrac{\sqrt3}{2}$}
\opt{C}{$\dfrac12$}
\opt{D}{$\dfrac34$}
\end{qoptions}

\qn{10} An object is placed at a distance of 15 cm from a convex lens of focal length 10 cm. On the other side of the lens, an convex mirror is placed at its focus such that the image formed by the combination coincides with the object itself. The radius of the convex mirror is
\begin{qoptions}
\opt{A}{20 cm}
\opt{B}{10 cm}
\opt{C}{15 cm}
\opt{D}{30 cm}
\end{qoptions}

\qn{11} In a circuit ideal diodes are given, find current drawn from battery?
\qfig{p1_phy_s1_q11_1}
\begin{qoptions}
\opt{A}{2.00 ampere}
\opt{B}{2.31 ampere}
\opt{C}{1.33 ampere}
\opt{D}{1.71 ampere}
\end{qoptions}

\qn{12} A star initially has $10^{40}$ deuterons. It produces nucleus via the processes\\
${}_1\mathrm{H}^2+{}_1\mathrm{H}^2\rightarrow{}_1\mathrm{H}^3+\mathrm{p}$\\
and ${}_1\mathrm{H}^2+{}_1\mathrm{H}^3\rightarrow{}_2\mathrm{He}^4+\mathrm{n}$\\
where the masses of the nuclei are :\\
$m({}_1\mathrm{H}^2)=2.014$ amu, $m(\mathrm p)=1.007$ amu, $m(\mathrm n)=1.008$ amu and $m({}_2\mathrm{He}^4)=4.001$ amu. If the average power radiated by the star is $10^{16}$ watt, the deuteron supply of the star is exhausted in a time of the order of :
\begin{qoptions}
\opt{A}{$10^6$ s}
\opt{B}{$10^8$ s}
\opt{C}{$10^{12}$ s}
\opt{D}{None}
\end{qoptions}

\qn{13} The radiation force experienced by body exposed to radiation of intensity $I$, assuming surface of body to be perfectly absorbing is :
\qfig{p1_phy_s1_q13_1}
\begin{qoptions}
\opt{A}{$\dfrac{\pi R^2I}{c}$}
\opt{B}{$\dfrac{\pi RHI}{c}$}
\opt{C}{$\dfrac{IRH}{2c}$}
\opt{D}{$\dfrac{IRH}{c}$}
\end{qoptions}

\qn{14} The electric potential between a proton and an electron is given by $V=V_0\ln\dfrac{r}{r_0}$, where $r_0$ is a constant. Assuming Bohr's model to be applicable, write variation of $r_n$ with $n$, $n$ being the principal quantum number :-
\begin{qoptions}
\opt{A}{$r_n\propto n$}
\opt{B}{$r_n\propto\dfrac1n$}
\opt{C}{$r_n\propto n^2$}
\opt{D}{$r_n\propto\dfrac1{n^2}$}
\end{qoptions}

\qn{15} The velocity $V$ of waves produced in water depends on their wavelength $\lambda$, the density of water $\rho$, and acceleration due to gravity $g$. The square of velocity is proportional to :-
\begin{qoptions}
\opt{A}{$\lambda^{-1}g^{-1}\rho^{-1}$}
\opt{B}{$\lambda g$}
\opt{C}{$\lambda\rho g$}
\opt{D}{$\lambda^2g^{-2}\rho^{-1}$}
\end{qoptions}

\qn{16} These blocks A, B and C of equal mass $m$ are placed one over the other on a smooth horizontal ground as shown in figure. Coefficient of friction between any two blocks of A, B and C is 1/2. The maximum value of mass of block D so that the blocks A, B and C move without slipping over each other is :-
\qfig{p1_phy_s1_q16_1}
\begin{qoptions}
\opt{A}{6m}
\opt{B}{5m}
\opt{C}{3m}
\opt{D}{4m}
\end{qoptions}

\qn{17} A particle moving along x-axis has acceleration $f$ at time $t$, given by : $f=f_0\left(1-\dfrac tT\right)$, where $f_0$ and $T$ are constants. The particle at $t=0$ has zero velocity. In the time interval between $t=0$ and the instant when $f=0$, the particle's velocity ($v_x$) is :--
\begin{qoptions}
\opt{A}{$1/2\,f_0T^2$}
\opt{B}{$f_0T^2$ (initial velocity is zero)}
\opt{C}{$1/2\,f_0T$}
\opt{D}{$f_0T$}
\end{qoptions}

\qn{18} A small disc of radius 2 cm is cut from a disc of radius 6 cm. If the distance between their centres is 3.2 cm, what is the shift in the centre of mass of the disc?
\begin{qoptions}
\opt{A}{0.4 cm}
\opt{B}{2.4 cm}
\opt{C}{1.8 cm}
\opt{D}{1.2 cm}
\end{qoptions}

\qn{19} A ball of mass $m$ moving horizontally at a speed $u$ collides with the bob of a simple pendulum at rest. The mass of the bob is $\dfrac m2$. If the collision is perfectly inelastic the height to which the composite ball rise after the collision is :-
\begin{qoptions}
\opt{A}{$\dfrac{u^2}{2g}$}
\opt{B}{$\dfrac{u^2}{8g}$}
\opt{C}{$\dfrac{2u^2}{9g}$}
\opt{D}{$\dfrac{u^2}{16g}$}
\end{qoptions}

\qn{20} A solid cylinder starts rolling from a height $h$ on an inclined plane. At same instant $t$, the ratio of its rotational K.E. and the total K.E. would be :-
\begin{qoptions}
\opt{A}{1 : 2}
\opt{B}{1 : 3}
\opt{C}{2 : 3}
\opt{D}{1 : 1}
\end{qoptions}

\sect{Section-II}

\qn{1} A Wheatstone's bridge is balanced with a resistance of 625 $\Omega$ in the third arm, where P, Q and S are in the 1$^{\text{st}}$, 2$^{\text{nd}}$ and 4$^{\text{th}}$ arm respectively. If P and Q are interchanged, the resistance in the third arm has to be increased by 51$\Omega$ to secure balance. The unknown resistance in the fourth arm is $6.5\times10^{N}\ \Omega$, then N will be :-
\par\nobreak
\qfig{p1_phy_s2_q01_1}

\qn{2} Two short electric dipoles of moment P and 64 P are placed in opposite direction on a line at a distance of 25 cm. The electric field will be zero at point between the dipoles whose distance from the dipole of moment P is (in cm) --

\qn{3} 2 kg of ice at $-20^\circ$C is mixed with 5 kg of water at 20$^\circ$C in an insulating vessel having a negligible heat capacity. Calculate the final mass of water (In Kg) remaining in the container. It is given that the specific heats of water and ice are 1 kcal/kg per $^\circ$C and 0.5 kcal/kg/$^\circ$C and the latent heat of fusion of ice is 80 kcal/kg :

\qn{4} Six identical rods are arranged as shown in the figure. The temperature of the junction B (In $^\circ$C) will be :-
\par\nobreak
\qfig{p1_phy_s2_q04_1}

\qn{5} A convex lens of focal length 20 cm and a concave lens of focal length $f$ are mounted coaxially 5 cm apart. Parallel beam of light incident on the convex lens emerges from the concave lens as a parallel beam. Then magnitude of `$f$' in cm is

\qn{6} A glass tube of length 1.0m is completely filled with water. A vibrating tuning fork of frequency 500 Hz is kept over the mouth of the tube and the water is drained out slowly at the bottom of the tube. If velocity of sound in air is 330 ms$^{-1}$, then the total number of resonances that occur will be :-

\qn{7} When a metallic surface is illuminated with radiation of wavelength $\lambda$, the stopping potential is $V$. If the same surface is illuminated with radiation of wavelength $2\lambda$, the stopping potential is $\dfrac V4$. The threshold wavelength for the metallic surface is :-

\qn{8} In the given circuit, the diode has a forward resistance of 8 $\Omega$ and infinite backward resistance. The value of $R$ is :-
\par\nobreak
\qfig{p1_phy_s2_q08_1}

\qn{9} The angular velocity of a body is $\vec\omega=2\hat i+3\hat j+4\hat k$ and a torque $\vec\tau=\hat i+2\hat j+3\hat k$ acts on it. The rotational power (in watt) will be :-

\qn{10} A simple pendulum of length 1 m is oscillating with an angular frequency 10 rad/s. The support of the pendulum starts oscillating up and down with a small angular frequency of 1 rad/s and an amplitude of $10^{-2}$ m. The relative change in the angular frequency of the pendulum is $\dfrac{n}{1000}$. Find value of n.

\partsec{2}{CHEMISTRY}
\sect{Section-I}

\qn{1} Which represents correct order of acidic strength ?
\begin{qoptions}
\opt{A}{NH$_3$ $>$ PH$_3$ $>$ AsH$_3$ $>$ SbH$_3$ $>$ BiH$_3$}
\opt{B}{La(OH)$_3$ $>$ Nd(OH)$_3$ $>$ Lu (OH)$_3$}
\opt{C}{NaOH $<$ KOH $<$ RbOH $<$ CsOH}
\opt{D}{CH$_4$ $<$ H$_2$S $<$ HCl $<$ HI}
\end{qoptions}

\qn{2} \textbf{Statement}-1 : Pb$^{4+}$ compounds are stronger oxidizing agents than Sn$^{4+}$ compounds.\\
\textbf{Statement}-2 : The higher oxidation states for the group 14 elements are more stable for the heavier members of the group due to inert pair effect.
\begin{qoptions}
\opt{A}{Statement-1 is true, Statement-2 is True, Statement-2 is a correct explanation for statement-1}
\opt{B}{Statement-1 is true, Statement-2 is True, Statement-2 is not a correct explanation for statement-1.}
\opt{C}{Statement-1 is True, Statement-2 is False.}
\opt{D}{Statement-1 is False, Statement-2 is True.}
\end{qoptions}

\qn{3} Which element is having lowest melting and boiling point?
\begin{qoptions}
\opt{A}{Ti}
\opt{B}{Cu}
\opt{C}{Zn}
\opt{D}{Mn}
\end{qoptions}

\qn{4} Which of the following is not formed when H$_2$S reacts with acidic K$_2$Cr$_2$O$_7$ solution ?
\begin{qoptions}
\opt{A}{K$_2$SO$_4$}
\opt{B}{CrSO$_4$}
\opt{C}{S}
\opt{D}{Cr$_2$(SO$_4$)$_3$}
\end{qoptions}

\qn{5} The number of unpaired electrons calculated in [Co(NH$_3$)$_6$]$^{3+}$ and [CoF$_6$]$^{3-}$ are :-
\begin{qoptions}
\opt{A}{4 and 4}
\opt{B}{0 and 2}
\opt{C}{2 and 4}
\opt{D}{0 and 4}
\end{qoptions}

\qn{6} Which of the following complex shows optical activity ?
\begin{qoptions}
\opt{A}{trans-[Co(ox)$_2$Cl$_2$]$^{3-}$}
\opt{B}{[Co(NH$_3$)$_3$Cl$_3$]}
\opt{C}{[Fe(edta)]$^-$}
\opt{D}{cis-[Pt(NH$_3$)$_2$Cl$_2$]}
\end{qoptions}

\qn{7} A diatomic molecule has dipole moment of 1.2 D. If its bond length is equal to $10^{-10}$ m then the fraction of an electronic charge on each atom will be :-
\begin{qoptions}
\opt{A}{42\%}
\opt{B}{52\%}
\opt{C}{37\%}
\opt{D}{25\%}
\end{qoptions}

\qn{8} (1) CO$_{(g)}$ + H$_2$O$_{(g)}$ $\rightleftharpoons$ CO$_{2(g)}$ + H$_{2(g)}$ ; K$_1$\\
(2) CH$_{4(g)}$ + H$_2$O$_{(g)}$ $\rightleftharpoons$ CO$_{(g)}$ + 3H$_{2(g)}$ ; K$_2$\\
(3) CH$_{4(g)}$ + 2H$_2$O$_{(g)}$ $\rightleftharpoons$ CO$_{2(g)}$ + 4H$_{2(g)}$ ; K$_3$\\
If the relation between K$_1$, K$_2$ and K$_3$ is $\mathrm{K_3}=\dfrac{(\mathrm{K_1})^x}{(\mathrm{K_2})^y}$ then calculate $(x-y)$
\begin{qoptions}
\opt{A}{2}
\opt{B}{3}
\opt{C}{1}
\opt{D}{0.5}
\end{qoptions}

\qn{9} Half life of a first order reaction :\\
A(aq.) $\rightarrow$ 2B(aq.) is 40 min at 27$^\circ$C and 20 min at 37$^\circ$C. The activation energy of reaction is ($\ln 2 = 0.7$, $R = 2.0$ cal mol$^{-1}$K$^{-1}$)
\begin{qoptions}
\opt{A}{19530 cal/ mol}
\opt{B}{13020 cal/ mol}
\opt{C}{26040 cal/ mol}
\opt{D}{6510 cal/ mol}
\end{qoptions}

\qn{10} Volume of water which must be added to 8M NaOH to make 2 lit. of 20 \% w/v NaOH solution is :
\begin{qoptions}
\opt{A}{1.25 lit.}
\opt{B}{0.75 lit.}
\opt{C}{1 lit.}
\opt{D}{None of these}
\end{qoptions}

\qn{11} Given, limiting equivalent conductance of BaCl$_2$ = $x_1$, H$_2$SO$_4$ = $x_2$, HCl = $x_3$\\
Specific conductance of the saturated BaSO$_4$ solution = $y$ Scm$^{-1}$. Thus K$_{\mathrm{SP}}$ of BaSO$_4$ is --
\begin{qoptions}
\opt{A}{$10^3y\,/\,2(x_1+x_2-2x_3)$}
\opt{B}{$10^6y^2\,/\,(x_1+x_2-2x_3)^2$}
\opt{C}{$10^6y^2\,/\,4(x_1+x_2-x_3)^2$}
\opt{D}{$(x_1+x^2-2x^3)\,/\,10^6y^2$}
\end{qoptions}

\qn{12} Two solutions A and B, each of 100 L was made by dissolving 4g of NaOH and 9.8 g of H$_2$SO$_4$ in water, respectively. The pH of the resultant solution obtained from mixing 40 L of solution A and 10 L of solution B is\underline{\hspace{3em}}.
\begin{qoptions}
\opt{A}{10}
\opt{B}{11}
\opt{C}{7}
\opt{D}{12}
\end{qoptions}

\qn{13} A hydrate of magnesium iodide has a formula MgI$_2$.xH$_2$O. If 1.055 g sample is heated to a constant weight of 0.695 g. What is the value of x ?
\begin{qoptions}
\opt{A}{4}
\opt{B}{2}
\opt{C}{6}
\opt{D}{8}
\end{qoptions}

\qn{14} What is the final temperature (in kelvin) of 0.10 mole monoatomic ideal gas that performs 75 cal of work adiabatically if the initial temperature is 227$^\circ$C ? (use $R = 2$ cal/K-mol)
\begin{qoptions}
\opt{A}{250}
\opt{B}{500}
\opt{C}{200}
\opt{D}{400}
\end{qoptions}

\qn{15} Among the following most stable resonating structure is :-
\begin{qoptions}
\opt{A}{\qfigi{p1_chem_s1_q15_1}}
\opt{B}{\qfigi{p1_chem_s1_q15_2}}
\opt{C}{\qfigi{p1_chem_s1_q15_5}}
\opt{D}{\qfigi{p1_chem_s1_q15_3}}
\end{qoptions}

\qn{16}\par\nobreak
\qfig{p1_chem_s1_q15_4}
A and B are respectively :-
\begin{qoptions}
\opt{A}{\qfigi{p1_chem_s1_q16_1} both}
\opt{B}{\qfigi{p1_chem_s1_q16_2}}
\opt{C}{\qfigi{p1_chem_s1_q16_3}}
\opt{D}{\qfigi{p1_chem_s1_q16_4}}
\end{qoptions}

\qn{17} Match the vitamins given in Column-I with the deficiency diseases listed in Column-II
\begin{center}
\begin{tabular}{|c|l|c|c|}
\hline
\multicolumn{2}{|c|}{\textbf{Column-I}} & \multicolumn{2}{c|}{\textbf{Column-II}}\\
\hline
(a) & Vitamin A & (p) & Osteomalacia\\
\hline
(b) & Vitamin C & (q) & Beri beri\\
\hline
(c) & Vitamin D & (r) & Scurvy\\
\hline
(d) & Vitamin B$_1$ & (s) & Xerophthalmia\\
\hline
\end{tabular}
\end{center}
\begin{qoptions}
\opt{A}{a $\rightarrow$ p; b $\rightarrow$ q; c $\rightarrow$ s; d $\rightarrow$ r}
\opt{B}{a $\rightarrow$ s; b $\rightarrow$ r; c $\rightarrow$ p; d $\rightarrow$ q}
\opt{C}{a $\rightarrow$ s; b $\rightarrow$ p; c $\rightarrow$ r; d $\rightarrow$ q}
\opt{D}{a $\rightarrow$ p; b $\rightarrow$ r; c $\rightarrow$ q; d $\rightarrow$ s}
\end{qoptions}

\qn{18}\par\nobreak
\qfig{p1_chem_s1_q18_1}
S is:
\begin{qoptions}
\opt{A}{\qfigi{p1_chem_s1_q18_2}}
\opt{B}{\qfigi{p1_chem_s1_q18_3}}
\opt{C}{\qfigi{p1_chem_s1_q18_4}}
\opt{D}{\qfigi{p1_chem_s1_q18_5}}
\end{qoptions}

\qn{19}\par\nobreak
\qfig{p1_chem_s1_q19_1}
Product (C) is:
\begin{qoptions}
\opt{A}{\qfigi{p1_chem_s1_q19_2}}
\opt{B}{\qfigi{p1_chem_s1_q19_3}}
\opt{C}{\qfigi{p1_chem_s1_q19_4}}
\opt{D}{\qfigi{p1_chem_s1_q19_5}}
\end{qoptions}

\qn{20} Which of the following statement is not correct ?
\begin{qoptions}
\opt{A}{Compound Ph$-$CH$_2-$Cl can undergo S$_\mathrm N$1 \& S$_\mathrm N$2 reactions}
\opt{B}{For compound $\mathrm{R-\overset{\overset{O}{\|}}{C}-CH_2-Cl}$ \ Rate of S$_\mathrm N$2 $<$ Rate of S$_\mathrm N$1}
\opt{C}{Rate of $<$ S$_\mathrm N$1 \qfigi{p1_chem_s1_q20_2}}
\opt{D}{Compound CH$_3$ $-$Cl will prefer to undergo S$_\mathrm N$2 when reacting with aq. KOH}
\end{qoptions}

\sect{Section-II}

\qn{1} The value of $(n\times\ell\times m)$ for the last electron entered in K is

\qn{2} How many of these molecules are non-polar ?\\
O$_3$, SO$_3$, PCl$_5$, CCl$_4$, SF$_4$, H$_2$O, PH$_3$, PCl$_3$F$_2$, XeF$_4$

\qn{3} How many of following orders are correct?\\
(a) Order of IE : O$^{-2}$ $>$ O$^+$ $>$ O $>$ O$^-$\\
(b) Order of EN : Zn $<$ Cd $<$ Hg\\
(c) Order of IE : B $>$ Tl $>$ Ga $>$ Al $>$ In\\
(d) Order of EA : S $>$ Se $>$ Te $>$ O\\
(e) Order of atomic radius : Ni $<$ Cu $<$ Zn\\
(f) Order of ionic radius : Na$^+$$>$Mg$^{+2}$$>$Li$^+$$>$Be$^{+2}$\\
(g) Order of IE$_2$ : N $<$ O $<$ F $<$ Ne

\qn{4} The half-life of a first order reaction is one hour. In what time will the concentration be reduced to $\dfrac78$ of its initial value :- ($\log 7 = 0.84$)

\qn{5} For HCl solution at 25$^\circ$C, Equivalent conductance at infinity dilution is 425 ohm$^{-1}$ cm$^2$ equivalent. The specific conductance of solution of HCl is 3.825 ohm$^{-1}$ cm$^{-1}$. If the degree of dissociation is 90\% the normality of the solution is :-

\qn{6} The solubility product of Cr(OH)$_3$ at 298 K is $6.0\times10^{-31}$. The concentration of hydroxide ions in a saturated solution of Cr(OH)$_3$ will be $(x\times10^{-31})^{1/4}$, then x is :

\qn{7} If x gram of non-volatile, non-electrolytic solute is dissolved in 114 g octane to reduce its vapour pressure to 80\% of that of pure octane then x is \underline{\hspace{2em}}\\
(Molar mass of solute = 40 g/mole)

\qn{8} How many chiral centres are present in
\par\nobreak
\qfig{p1_chem_s2_q08_1} Fluoromethanolone

\qn{9} How many of the following give positive isocyanide test\\
(a) \qfigi{p1_chem_s2_q09_1}\quad (b) \qfigi{p1_chem_s2_q09_2}\quad (c) Ph$-$NH$_2$\\
(d) \qfigi{p1_chem_s2_q09_3}\quad (e) \qfigi{p1_chem_s2_q09_4}\quad (f) \qfigi{p1_chem_s2_q09_7}\\
(g) \qfigi{p1_chem_s2_q09_6}\quad (h) \qfigi{p1_chem_s2_q09_5}

\qn{10} How many species is/are $-$M/$-$R groups ?\\
(a) $-$OH\quad (b) $-$NO$_2$\quad (c) $-$CN\\
(d) $-$CH$_3$\quad (e) $-$O$^{\ominus}$\quad (f) $-\overset{\oplus}{\mathrm N}$H$_3$\\
(g) \qfigi{p1_chem_s2_q10_1}\quad (h) $-$CH=NH

\partsec{3}{MATHEMATICS}
\sect{Section-I}

\qn{1} The integral $\displaystyle\int\dfrac{(x^8-x^2)dx}{(x^{12}+3x^6+1)\tan^{-1}\left(x^3+\frac1{x^3}\right)}$ is equal to :
\begin{qoptions}
\opt{A}{$\log_e\left(\left|\tan^{-1}\left(x^3+\dfrac1{x^3}\right)\right|\right)^{1/3}+C$}
\opt{B}{$\log_e\left(\left|\tan^{-1}\left(x^3+\dfrac1{x^3}\right)\right|\right)^{1/2}+C$}
\opt{C}{$\log_e\left(\left|\tan^{-1}\left(x^3+\dfrac1{x^3}\right)\right|\right)+C$}
\opt{D}{$\log_e\left(\left|\tan^{-1}\left(x^3+\dfrac1{x^3}\right)\right|\right)^{3}+C$}
\end{qoptions}

\qn{2} The area enclosed by the curves $xy+4y=16$ and $x+y=6$ is equal to :
\begin{qoptions}
\opt{A}{$28-30\log_e2$}
\opt{B}{$30-28\log_e2$}
\opt{C}{$30-32\log_e2$}
\opt{D}{$32-30\log_e2$}
\end{qoptions}

\qn{3} Let $A=\begin{bmatrix}2&1&2\\6&2&11\\3&3&2\end{bmatrix}$ and $P=\begin{bmatrix}1&2&0\\5&0&2\\7&1&5\end{bmatrix}$. The sum of the prime factors of $\left|P^{-1}AP-2I\right|$ is equal to
\begin{qoptions}
\opt{A}{26}
\opt{B}{27}
\opt{C}{66}
\opt{D}{23}
\end{qoptions}

\qn{4} An integer is chosen at random from the integers 1, 2, 3$\ldots$50. The probability that the chosen integer is a multiple of atleast one of 4, 6 and 7 is :
\begin{qoptions}
\opt{A}{$\dfrac{8}{25}$}
\opt{B}{$\dfrac{21}{50}$}
\opt{C}{$\dfrac{9}{50}$}
\opt{D}{$\dfrac{14}{25}$}
\end{qoptions}

\qn{5} If the system of linear equations\\
$x-2y+z=-4$\\
$2x+\alpha y+3z=5$\\
$3x-y+\beta z=3$\\
has infinitely many solutions, then $12\alpha+13\beta$ is equal to
\begin{qoptions}
\opt{A}{60}
\opt{B}{64}
\opt{C}{54}
\opt{D}{58}
\end{qoptions}

\qn{6} If $\log_ea$, $\log_eb$, $\log_ec$ are in an A.P. and $\log_ea-\log_e2b$, $\log_e2b-\log_e3c$, $\log_e3c-\log_ea$ are also in an A.P., then $a:b:c$ is equal to
\begin{qoptions}
\opt{A}{9 : 6 : 4}
\opt{B}{16 : 4 : 1}
\opt{C}{25 : 10 : 4}
\opt{D}{6 : 3 : 2}
\end{qoptions}

\qn{7} Let Ajay will not appear in JEE exam with probability $p=\dfrac27$, while both Ajay and Vijay will appear in the exam with probability $q=\dfrac15$. Then the probability, that Ajay will appear in the exam and Vijay will not appear is :
\begin{qoptions}
\opt{A}{$\dfrac{9}{35}$}
\opt{B}{$\dfrac{18}{35}$}
\opt{C}{$\dfrac{24}{35}$}
\opt{D}{$\dfrac{3}{35}$}
\end{qoptions}

\qn{8} If $A=\begin{bmatrix}\sqrt2&1\\-1&\sqrt2\end{bmatrix}$, $B=\begin{bmatrix}1&0\\1&1\end{bmatrix}$, $C=ABA^T$ and $X=A^TC^2A$, then $\det X$ is equal to :
\begin{qoptions}
\opt{A}{243}
\opt{B}{729}
\opt{C}{27}
\opt{D}{891}
\end{qoptions}

\qn{9} If $\alpha,\beta$ are the roots of the equation, $x^2-x-1=0$ and $S_n=2023\alpha^n+2024\beta^n$, then
\begin{qoptions}
\opt{A}{$2S_{12}=S_{11}+S_{10}$}
\opt{B}{$S_{12}=S_{11}+S_{10}$}
\opt{C}{$2S_{11}=S_{12}+S_{10}$}
\opt{D}{$S_{11}=S_{10}+S_{12}$}
\end{qoptions}

\qn{10} Function $f(x)=\dfrac{\lambda\sin x+3\cos x}{2\sin x+6\cos x}$ is monotonic increasing when :-
\begin{qoptions}
\opt{A}{$\lambda<1$}
\opt{B}{$\lambda>1$}
\opt{C}{$\lambda<2$}
\opt{D}{$\lambda>2$}
\end{qoptions}

\qn{11} If\\\resizebox{\linewidth}{!}{$f(x)=\begin{cases}\dfrac{1-\sin x}{(\pi-2x)^2}\cdot\dfrac{\log\sin x}{\log\left(1+\pi^2-4\pi x+4x^2\right)} & ,\ x\ne\dfrac\pi2\\[6pt] k & ,\ x=\dfrac\pi2\end{cases}$}\\
is continuous at $x=\pi/2$, then $k$ is -
\begin{qoptions}
\opt{A}{$-\dfrac1{16}$}
\opt{B}{$-\dfrac1{32}$}
\opt{C}{$-\dfrac1{64}$}
\opt{D}{$-\dfrac1{128}$}
\end{qoptions}

\qn{12} A possible value of $\tan\left(\dfrac14\sin^{-1}\dfrac{\sqrt{63}}{8}\right)$ is :
\begin{qoptions}
\opt{A}{$\dfrac1{\sqrt7}$}
\opt{B}{$2\sqrt2-1$}
\opt{C}{$\sqrt7-1$}
\opt{D}{$\dfrac1{2\sqrt2}$}
\end{qoptions}

\qn{13} If $x\in(\pi,2\pi)$ \& $\dfrac{\sqrt{1+\cos x}+\sqrt{1-\cos x}}{\sqrt{1+\cos x}-\sqrt{1-\cos x}}=\cot\left(a+\dfrac x2\right)$ then `$a$' equals :-
\begin{qoptions}
\opt{A}{$\dfrac\pi4$}
\opt{B}{$\dfrac\pi2$}
\opt{C}{$\dfrac\pi3$}
\opt{D}{None of these}
\end{qoptions}

\qn{14} If $\displaystyle\sum_{i=1}^{5}(x_i-10)=5$ and $\displaystyle\sum_{i=1}^{5}(x_i-10)^2=25$, then standard deviation of observations $2x_1+7$, $2x_2+7$, $2x_3+7$, $2x_4+7$ and $2x_5+7$ is equal to-
\begin{qoptions}
\opt{A}{8}
\opt{B}{16}
\opt{C}{4}
\opt{D}{2}
\end{qoptions}

\qn{15} $96\cos\dfrac\pi{33}\cos\dfrac{2\pi}{33}\cos\dfrac{4\pi}{33}\cos\dfrac{8\pi}{33}\cos\dfrac{16\pi}{33}$ is equal to
\begin{qoptions}
\opt{A}{3}
\opt{B}{2}
\opt{C}{4}
\opt{D}{1}
\end{qoptions}

\qn{16} Vector $\hat a$ in the plane of $\vec b=2\hat i+\hat j$ and $\vec c=\hat i-\hat j+\hat k$ is such that it is equally inclined to $\vec b$ and $\vec d$ where $\vec d=\hat j+2\hat k$. The value of $\hat a$ is :-
\begin{qoptions}
\opt{A}{$\dfrac{\hat i+\hat j+\hat k}{\sqrt3}$}
\opt{B}{$\dfrac{\hat i-\hat j+\hat k}{\sqrt3}$}
\opt{C}{$\dfrac{2\hat i+\hat j}{\sqrt5}$}
\opt{D}{None of these}
\end{qoptions}

\qn{17} Given $A(1,2)$ \& $B(3,7)$ are two vertices of the $\Delta ABC$. If the locus of centroid of $\Delta ABC$ is $2x-y=0$ then the minimum distance between the locus of the vertex $C$ from the line $2x-y=0$ is-
\begin{qoptions}
\opt{A}{0}
\opt{B}{1}
\opt{C}{$\dfrac2{\sqrt5}$}
\opt{D}{$\dfrac1{\sqrt5}$}
\end{qoptions}

\qn{18} On portion of straight line $x+y=2$ which is intercepted between the axes, a square away from the origin is constructed. $p$ denotes the perpendicular distance of a side of square from origin, then maximum value of $p$ is :
\begin{qoptions}
\opt{A}{$\sqrt2$}
\opt{B}{$2\sqrt2$}
\opt{C}{$3\sqrt2$}
\opt{D}{$4\sqrt2$}
\end{qoptions}

\qn{19} The co-ordinates of a point on the parabola $y^2=8x$ whose focal distance is 4 is :-
\begin{qoptions}
\opt{A}{$(2,\pm4)$}
\opt{B}{$(\pm2,4)$}
\opt{C}{$(-2,\pm4)$}
\opt{D}{$(\pm2,-4)$}
\end{qoptions}

\qn{20} The straight lines whose direction cosines are given by $a\ell+bm+cn=0$, $fmn+gn\ell+h\ell m=0$ are perpendicular, if :-
\begin{qoptions}
\opt{A}{$\dfrac fa+\dfrac gb+\dfrac hc=0$}
\opt{B}{$\sqrt{\dfrac af}+\sqrt{\dfrac bg}+\sqrt{\dfrac ch}=0$}
\opt{C}{$\sqrt{af}=\sqrt{bg}=\sqrt{ch}$}
\opt{D}{$\sqrt{\dfrac af}=\sqrt{\dfrac bg}=\sqrt{\dfrac ch}$}
\end{qoptions}

\sect{Section-II}

\qn{1} If the solution curve, of the differential equation $\dfrac{dy}{dx}=\dfrac{x+y-2}{x-y}$ passing through the point $(2,1)$ is $\tan^{-1}\left(\dfrac{y-1}{x-1}\right)-\dfrac1\beta\log_e\left(\alpha+\left(\dfrac{y-1}{x-1}\right)^2\right)=\log_e|x-1|$, then $5\beta+\alpha$ is equal to

\qn{2} If $x=x(t)$ is the solution of the differential equation $(t+1)dx=\left(2x+(t+1)^4\right)dt$, $x(0)=2$, then, $x(1)$ equals \underline{\hspace{4em}}.

\qn{3} If $\alpha$ satisfies the equation $x^2+x+1=0$ and $(1+\alpha)^7=A+B\alpha+C\alpha^2$, $A,B,C\ge0$, then $5(3A-2B-C)$ is equal to \underline{\hspace{3em}}.

\qn{4} In the expansion of $(1+x)(1-x^2)\left(1+\dfrac3x+\dfrac3{x^2}+\dfrac1{x^3}\right)^5$, $x\ne0$, the sum of the coefficient of $x^3$ and $x^{-13}$ is equal to \underline{\hspace{3em}}

\qn{5} The points of extremum of the function $F(x)=\displaystyle\int_1^{x}e^{-t^2/2}(1-t^2)\,dt$ are :-

\qn{6} If $f(x)$ is a polynomial satisfying $f(x)f(1/x)=f(x)+f(1/x)$ and $f(2)>1$ then $\displaystyle\lim_{x\to1}f(x)$ is :-

\qn{7} If $\sin^2\alpha+\sin^4\alpha=1$ \& $\cos^8\alpha-\sin^8\alpha=\left(b-9\sqrt a\right)$, then $\dfrac ba$ is equal to

\qn{8} Total number of equivalence relations defined in the set $S=\{a,b,c\}$ is

\qn{9} If the shortest distance between the line $\vec r=\left(-\hat i+3\hat k\right)+\lambda\left(\hat i-a\hat j\right)$ and $\vec r=\left(-\hat j+2\hat k\right)+\mu\left(\hat i-\hat j+\hat k\right)$ is $\sqrt{\dfrac23}$, then the integral value of $a$ is equal to

\qn{10} If the area of an equilateral triangle inscribed in the circle, $x^2+y^2+10x+12y+c=0$ is $27\sqrt3$ sq. units then $c$ is equal to
\end{multicols}
\end{document}
